% Section 1 -- drafted by Claude (2026-10-21), revised after review.
\section{Introduction}\label{sec:intro}

The Predictive AI Evaluation Competition \citep{truong2026paiec} scores predictors on benchmarks they have never seen. For an AI system (a subject) and an item of a hidden benchmark, the predictor gives the probability that the subject succeeds, without running it. It starts from the subject's recorded fields and the item's text. Codabench scores each upload on a fresh random draw of subject--benchmark pairs (8 or 9 in each draw I saw) as formative feedback. At budgets of 1 to 31 labels the predictor also sees the subject's outcomes on other items of that benchmark, with the labels of other subjects in the same draw. The Brier ALC (area under the learning curve, Eq.~\eqref{eq:alc}) averages the Brier scores at six budgets from 0 to 31 labels and puts 70\% of its weight on budgets of 7 labels or fewer. This is a small instance of predictable AI in the sense of \citet{zhou2026predictable}.

Item response models have been applied to AI evaluation data \citep{lalor2016building,martinezplumed2019irt,rodriguez2021evaluation} and can estimate a model's score from a few items \citep{polo2024tinybenchmarks,truong2025reeval}; \citet{truong2025reeval} also predict a new item's difficulty from its content, with a predictor fit to responses on other items. Without such a predictor, these methods need responses on the same items, and at the start of a hidden benchmark there are none. \citet{zhang2025benchmark} found that a random sample with a regression adjustment beats most benchmark-prediction methods, and that when new models are more accurate than those seen before, none of these methods consistently beats a plain average over a random sample. Estimating difficulty from item text alone is an old problem in education \citep{benedetto2023survey}, and LLM judgments of the difficulty of human test items match measured difficulty only weakly to moderately \citep{chrzan2026consensus,li2026llms}.

I built the entry as a chain of small Bayesian models and decided most additions by a rule written before its test (Appendix~\ref{app:rules} lists the exceptions). V1 places a normal prior on a subject's ability, centered on a score from the subject's metadata, and updates it on a grid; V2 pools item difficulties from the labels of other subjects in the same draw. V3 keeps V2 but drops the metadata score from the prior, so that it predicts 0.5 before the first label. Most of the offline gain came from the two LLM inputs added last: V4 uses Claude Opus 5.5's rating of each item's difficulty as a prior, and V5, the final entry, adds the same model's guess of how often the subject succeeds on the benchmark. Both inputs are calibrated on public data and enter with a trust level chosen to limit the loss if they turn out to be wrong. Because Claude may have seen these public benchmarks in training, the offline gains may overstate what transfers to new ones (Section~\ref{sec:limitations}). The contributions are:
\begin{itemize}
  \item a local harness that simulates the scoring protocol and whose ALC computation reproduces all 68 per-pair scores I kept from Codabench, together with the rules fixed before each test and the ideas that failed (Sections~\ref{sec:data} and~\ref{sec:analysis}, Appendix~\ref{app:rules});
  \item two ways to use an LLM inside a Bayesian predictor, as a gated prior on item difficulty and as a prior on subject ability, which improve offline Brier ALC on rated items by 0.0144 (SE 0.0011) and a further 0.0135 (SE 0.0013) (Sections~\ref{sec:method} and~\ref{sec:results});
  \item a submission that calls an LLM API during hosted scoring despite hidden program output, no package installation and workers that restart at every budget (Section~\ref{sec:llm});
  \item an account of what single formative draws can show: metadata priors helped offline but hurt at budget 0 on all three hosted draws, and the final entry's budget-0 Brier fell to 0.180 from the 0.25 of a constant 0.5 (Section~\ref{sec:results}).
\end{itemize}
